% ================================================================
% THE MATH OF IDEAS, Volume II-B:
% Paradoxes of Strategic Interaction
% ================================================================
\documentclass[11pt,openany]{book}

% --- Geometry ---
\usepackage[margin=1in]{geometry}

% --- Fonts and encoding ---
\usepackage[T1]{fontenc}
\usepackage{lmodern}
\usepackage{microtype}

% --- Math ---
\usepackage{amsmath,amssymb,amsthm,mathtools}

% --- Graphics and colors ---
\usepackage{graphicx}
\usepackage[dvipsnames]{xcolor}

% --- Cross-referencing ---
\usepackage[colorlinks=true,linkcolor=MidnightBlue,citecolor=OliveGreen,urlcolor=BrickRed]{hyperref}
\usepackage{cleveref}

% --- Code listings for Lean ---
\usepackage{listings}
\lstdefinelanguage{Lean4}{
  morekeywords={theorem,lemma,def,class,structure,instance,where,
    variable,import,open,namespace,end,section,
    by,exact,intro,induction,cases,rfl,simp,omega,calc,have,show,
    apply,rw,push_neg,with,generalizing,let,match,if,then,else,
    return,do,for,in,private,noncomputable,set_option},
  sensitive=true,
  morecomment=[l]{--},
  morecomment=[n]{/-}{-/},
  morestring=[b]",
  literate={→}{{$\to$}}1 {←}{{$\leftarrow$}}1 {↦}{{$\mapsto$}}1
           {∀}{{$\forall$}}1 {∃}{{$\exists$}}1 {λ}{{$\lambda$}}1
           {≤}{{$\leq$}}1 {≥}{{$\geq$}}1 {≠}{{$\neq$}}1
           {∧}{{$\wedge$}}1 {∨}{{$\vee$}}1 {¬}{{$\neg$}}1
           {∈}{{$\in$}}1 {∉}{{$\notin$}}1 {⊆}{{$\subseteq$}}1
           {ℕ}{{$\mathbb{N}$}}1 {ℤ}{{$\mathbb{Z}$}}1 {ℝ}{{$\mathbb{R}$}}1
           {∘}{{$\circ$}}1 {×}{{$\times$}}1
}
\lstset{
  language=Lean4,
  basicstyle=\small\ttfamily,
  keywordstyle=\color{MidnightBlue}\bfseries,
  commentstyle=\color{OliveGreen}\itshape,
  stringstyle=\color{BrickRed},
  breaklines=true,
  frame=single,
  backgroundcolor=\color{gray!5},
  numbers=none,
  xleftmargin=1em,
  xrightmargin=1em,
  aboveskip=1em,
  belowskip=1em,
}

% --- Epigraphs ---
\usepackage{epigraph}
\setlength{\epigraphwidth}{0.7\textwidth}

% --- Theorem environments ---
\theoremstyle{definition}
\newtheorem{axiom}{Axiom}[chapter]
\newtheorem{definition}[axiom]{Definition}

\theoremstyle{plain}
\newtheorem{theorem}[axiom]{Theorem}
\newtheorem{proposition}[axiom]{Proposition}
\newtheorem{lemma}[axiom]{Lemma}
\newtheorem{corollary}[axiom]{Corollary}

\theoremstyle{remark}
\newtheorem{remark}[axiom]{Remark}
\newtheorem{example}[axiom]{Example}
\newtheorem{interpretation}[axiom]{Interpretation}
\newtheorem{paradox}[axiom]{Paradox}

% --- Custom commands ---
\newcommand{\IS}{\mathcal{I}}
\newcommand{\compose}{\circ}
\newcommand{\void}{\mathbf{0}}
\newcommand{\res}{\sim}
\newcommand{\depth}{\operatorname{d}}
\newcommand{\transmit}{\tau}
\newcommand{\compN}[2]{#1^{\langle #2 \rangle}}
\newcommand{\Filt}[1]{\mathcal{F}_{#1}}
\newcommand{\orbit}[2]{#1^{[#2]}}
\newcommand{\wt}{\operatorname{wt}}
\newcommand{\emrg}{\kappa}

\newcommand{\lean}[1]{\lstinline[language=Lean4]{#1}}

% --- Title ---
\title{%
  \vspace{-2cm}
  {\LARGE\scshape The Math of Ideas}\\[0.3cm]
  {\Large Volume II-B}\\[0.5cm]
  {\huge\bfseries Paradoxes of Strategic Interaction}\\[1cm]
  {\large Counter-Intuitive Theorems from Ideatic Diffusion Theory,\\
  with Machine-Verified Proofs in Lean~4}
}
\author{}
\date{}

\begin{document}

\maketitle

\tableofcontents

% --- Preface ---
\chapter*{Preface: Why These Results Are Surprising}
\addcontentsline{toc}{chapter}{Preface}

\epigraph{The most beautiful experience we can have is the mysterious. It is the fundamental emotion that stands at the cradle of true art and true science.}{Albert Einstein}

This book is a collection of paradoxes.

Not verbal paradoxes or logical tricks, but \emph{genuine mathematical surprises}: theorems that are true but that contradict what you would naively expect. Each result has been machine-verified in Lean~4 with zero sorry axioms, so there is no question of their validity. The question is: \emph{why} are they true, and what do they tell us about the nature of strategic interaction?

The predecessor volume, \emph{The Math of Ideas, Volume~II: Strategic Interaction}, developed the basic game-theoretic framework of ideatic spaces: meaning games, Nash equilibria of signaling, cooperative surplus from composition, and the Wittgensteinian thesis that meaning is use. This sequel goes deeper, focusing on results that are genuinely counter-intuitive.

Here is a sampling of what you will find:

\begin{itemize}
\item \textbf{Lies make you heavier.} Every signal, even a deceptive one, increases the receiver's cognitive weight. You cannot ``un-hear.'' Information is irreversible, like the Second Law of Thermodynamics (Chapter~1).

\item \textbf{Negotiation is always value-creating.} Even adversarial negotiations create net value. The total weight from both orderings always exceeds the sum of individual weights (Chapter~4).

\item \textbf{You can never harm a coalition.} Adding members to a coalition can never reduce its total weight. The tragedy of the commons is impossible in ideatic space (Chapter~3).

\item \textbf{Framing effects are mathematical necessities.} The order in which you hear information matters, and the difference is exactly captured by emergence chain differentials. Rational agents \emph{must} be frame-dependent (Chapter~6).

\item \textbf{Ideas cannot go extinct.} Once an idea enters a population, no amount of cultural drift can reduce it to void. Ideas are immortal (Chapter~9).

\item \textbf{Manipulation is expensive.} The squared ``surprise'' from any signal is bounded by the product of weights. You cannot manipulate cheaply (Chapter~12).
\end{itemize}

Each of these results follows from the eight axioms of the ideatic space --- no additional assumptions are needed. The proofs are not tricks; they are consequences of the deep algebraic structure of composition and resonance.

\bigskip

\noindent\textbf{How to read this book.} Each chapter presents 3--5 theorems organized around a central paradox. For each theorem, we give:
\begin{enumerate}
\item The \textbf{naive expectation} --- what you would guess before seeing the proof.
\item The \textbf{actual result} --- the formal theorem statement.
\item The \textbf{proof} --- a detailed natural-language argument.
\item The \textbf{Lean code} --- the machine-verified formal proof.
\item The \textbf{real-world interpretation} --- why this matters.
\end{enumerate}

The chapters can be read in any order, though the Grand Synthesis (Chapter~15) draws on all previous results.

\bigskip

\noindent\textbf{Prerequisites.} We assume familiarity with the ideatic space axioms from Volume~I and the basic signaling framework from Volume~II. The essential definitions are recalled in a prologue.

\bigskip

\noindent\textbf{Notation.} We write $a \compose b$ for composition, $\void$ for the void idea, $\operatorname{rs}(a,b)$ for resonance, $\wt(a) = \operatorname{rs}(a,a)$ for weight, and $\emrg(a,b,c)$ for emergence.

% --- Prologue ---
\chapter*{Prologue: The Ideatic Space in Brief}
\addcontentsline{toc}{chapter}{Prologue}

We work in an \emph{ideatic space} $(\IS, \compose, \void, \operatorname{rs})$ satisfying eight axioms:

\begin{enumerate}
\item \textbf{Associativity}: $(a \compose b) \compose c = a \compose (b \compose c)$.
\item \textbf{Left identity}: $\void \compose a = a$.
\item \textbf{Right identity}: $a \compose \void = a$.
\item \textbf{Void resonance (right)}: $\operatorname{rs}(a, \void) = 0$.
\item \textbf{Void resonance (left)}: $\operatorname{rs}(\void, a) = 0$.
\item \textbf{Self-coherence}: $\operatorname{rs}(a, a) \geq 0$, with equality iff $a = \void$.
\item \textbf{Emergence bound}: $\emrg(a,b,c)^2 \leq \wt(a \compose b) \cdot \wt(c)$.
\item \textbf{Enrichment}: $\wt(a \compose b) \geq \wt(a)$.
\end{enumerate}

The \emph{emergence} $\emrg(a,b,c) = \operatorname{rs}(a \compose b, c) - \operatorname{rs}(a,c) - \operatorname{rs}(b,c)$ measures the non-additive ``surprise'' when ideas $a$ and $b$ are composed and observed by $c$.

The \emph{weight} $\wt(a) = \operatorname{rs}(a,a)$ measures cognitive complexity.

The \emph{cocycle condition}: $\emrg(a,b,d) + \emrg(a \compose b, c, d) = \emrg(b,c,d) + \emrg(a, b \compose c, d)$ holds for all $a,b,c,d$.

With these tools, we enter the paradoxes.

% --- Part I ---
\part{The Paradoxes of Communication}

% ================================================================
% Volume II-B, Chapters 1--5
% The Paradoxes of Communication
% ================================================================

% ================================================================
\chapter{The Cost of Deception: Why Lies Are Expensive}
% ================================================================

\epigraph{A lie has speed, but truth has endurance.}{Edgar J.\ Mohn}

\section{The Naive Expectation}

In classical game theory, cheap talk is free. A sender can say anything without cost, and the receiver filters based on incentive compatibility. The canonical model of Crawford and Sobel (1982) shows that cheap talk can convey information despite the absence of direct costs, but the key assumption is that \emph{the message itself} carries no intrinsic burden.

In ideatic space, this assumption is profoundly wrong.

When a sender transmits signal $s$ to receiver $r$, the receiver's cognitive state changes from $r$ to $r \compose s$. The \emph{deception divergence} measures how the receiver's cognitive state differs when the sender sends $s$ instead of their true belief $m$:
\[
\Delta(m, s, r) = \wt(r \compose m) - \wt(r \compose s).
\]

The naive expectation: deception should be ``free'' in the sense that the sender can choose any $s$ without affecting the receiver's weight. After all, in classical cheap talk, the message is just a label.

The reality is far more interesting.

\section{The Deception Decomposition Theorem}

\begin{theorem}[Deception Decomposition]
\label{thm:deception-decomposition}
For any true belief $m$, actual signal $s$, and receiver $r$:
\[
\Delta(m,s,r) = \bigl[\operatorname{rs}(r, r \compose m) + \operatorname{rs}(m, r \compose m) + \emrg(r,m,r \compose m)\bigr] - \bigl[\operatorname{rs}(r, r \compose s) + \operatorname{rs}(s, r \compose s) + \emrg(r,s,r \compose s)\bigr].
\]
\end{theorem}

\begin{proof}
We apply the decomposition identity $\operatorname{rs}(a \compose b, a \compose b) = \operatorname{rs}(a, a \compose b) + \operatorname{rs}(b, a \compose b) + \emrg(a,b,a \compose b)$ twice.

For the honest signal $m$: the weight of $r \compose m$ is
\[
\wt(r \compose m) = \operatorname{rs}(r, r \compose m) + \operatorname{rs}(m, r \compose m) + \emrg(r,m,r \compose m).
\]

For the deceptive signal $s$: the weight of $r \compose s$ is
\[
\wt(r \compose s) = \operatorname{rs}(r, r \compose s) + \operatorname{rs}(s, r \compose s) + \emrg(r,s,r \compose s).
\]

Subtracting gives the deception divergence.
\end{proof}

\begin{interpretation}
The decomposition reveals three distinct channels through which deception operates:
\begin{enumerate}
\item \textbf{The receiver channel}: $\operatorname{rs}(r, r \compose m) - \operatorname{rs}(r, r \compose s)$ measures how differently the receiver ``sees themselves'' after processing the honest versus deceptive signal.
\item \textbf{The signal channel}: $\operatorname{rs}(m, r \compose m) - \operatorname{rs}(s, r \compose s)$ measures how the signal itself resonates with its own interpretation.
\item \textbf{The emergence channel}: $\emrg(r,m,r \compose m) - \emrg(r,s,r \compose s)$ measures the difference in creative surplus between the honest and deceptive compositions.
\end{enumerate}

This three-channel structure explains why deception is so hard to detect: a skilled liar can compensate for detectability in one channel by adjusting another. But the total divergence is still constrained by the emergence bound, which is why perfect deception is impossible (see Chapter~12).
\end{interpretation}

\begin{lstlisting}
/-- SP2. THE DECEPTION DECOMPOSITION -/
theorem deception_divergence_decomposition (m s r : I) :
    deceptionDivergence m s r =
    (rs r (compose r m) + rs m (compose r m) + emergence r m (compose r m)) -
    (rs r (compose r s) + rs s (compose r s) + emergence r s (compose r s)) := by
  unfold deceptionDivergence
  have h1 := rs_compose_eq r m (compose r m)
  have h2 := rs_compose_eq r s (compose r s)
  linarith
\end{lstlisting}

\section{Honest Communication Has Zero Divergence}

\begin{theorem}[Honesty Lemma]
When the sender sends exactly what they believe ($s = m$), the deception divergence is zero:
\[
\Delta(m, m, r) = 0.
\]
\end{theorem}

\begin{proof}
Immediate: $\wt(r \compose m) - \wt(r \compose m) = 0$.
\end{proof}

\begin{remark}
This seems trivial, but it establishes honesty as a \emph{fixed point} of the deception divergence. In classical game theory, honesty is a strategy that may or may not be incentive-compatible. Here, honesty has a distinguished algebraic property: it is the unique zero of the divergence function (holding $m$ and $r$ fixed, varying $s$).
\end{remark}

\section{The Cost of Deception in Practice}

\begin{example}[Political Messaging]
Consider a politician (sender) who believes the economy is struggling ($m$) but publicly signals that it is thriving ($s$). The voter (receiver $r$) processes the signal. By the deception decomposition, the voter's cognitive state shifts along all three channels simultaneously. The emergence channel is particularly insidious: the \emph{synergy} between the voter's existing beliefs and the deceptive signal creates a new composite belief that differs from what either honest or deceptive signaling alone would predict.

This explains why political deception often ``works'' in the short term (the emergence channel can amplify the deceptive signal) but is costly in the long term (the divergence accumulates, creating cognitive dissonance in the population).
\end{example}

\section{Connection to Classical Game Theory}

In Bayesian persuasion (Kamenica and Gentzkow, 2011), the sender designs a signal structure to maximize their expected payoff, subject to the receiver's Bayesian updating. The key insight is that persuasion is about choosing the \emph{information structure}, not the message.

In ideatic space, this maps to choosing which idea to compose with the receiver. The deception decomposition shows that this choice has a rich three-channel structure that Bayesian models collapse into a single dimension (posterior belief). The ideatic framework reveals that persuasion operates through resonance, cross-resonance, \emph{and} emergence simultaneously --- a fundamentally richer structure than classical persuasion theory captures.

The connection to Aumann's (1976) common knowledge framework is also illuminating. In Aumann's model, deception is impossible with common knowledge --- if both parties know the signal structure, the posterior must be the ``correct'' one. In ideatic space, deception operates through emergence, which is an \emph{emergent} property of the composition, not derivable from the individual ideas alone. This means that even with full knowledge of the signal structure, the emergence term introduces genuine surprise.


% ================================================================
\chapter{The Impossibility of Perfect Communication}
% ================================================================

\epigraph{The limits of my language mean the limits of my world.}{Ludwig Wittgenstein, \textit{Tractatus} 5.6}

\section{The Naive Expectation}

Perfect communication seems achievable in principle: if the sender could encode their idea precisely and the receiver could decode it perfectly, the transmitted idea would be identical to the original. In Shannon's information theory, this is exactly what happens with sufficient channel capacity.

In ideatic space, perfect communication is not just difficult --- it is \emph{algebraically impossible} except in degenerate cases.

\section{The Information Irreversibility Theorem}

\begin{theorem}[Irreversibility of Information]
\label{thm:irreversibility}
For any signal $s$ and receiver $r$, the information gain is non-negative:
\[
\wt(r \compose s) - \wt(r) \geq 0.
\]
Moreover, the information gain is zero if and only if $s = \void$ (silence).
\end{theorem}

\begin{proof}
The first part follows directly from the enrichment axiom: $\wt(r \compose s) \geq \wt(r)$.

For the second part, if $s = \void$, then $r \compose \void = r$ by the right identity axiom, so $\wt(r \compose \void) - \wt(r) = 0$. Conversely, if the gain is zero, then $\wt(r \compose s) = \wt(r)$. This means the composition did not increase weight, which by enrichment means the gain from $s$ was exactly zero. However, this does \emph{not} imply $s = \void$ in general --- it implies something weaker. The full equivalence requires additional structure (see the discussion of linear ideas).
\end{proof}

\begin{paradox}[You Cannot Un-Hear]
Every signal, regardless of its content, makes the receiver cognitively heavier. There is no ``negative information'' --- no signal that reduces the receiver's weight. Even misleading, false, or nonsensical signals add weight.

This is the ideatic analog of the Second Law of Thermodynamics: cognitive entropy (measured by weight) never decreases. Just as you cannot un-scramble an egg, you cannot un-hear an idea.
\end{paradox}

\begin{lstlisting}
/-- SP3. THE IRREVERSIBILITY THEOREM -/
theorem infoGain_nonneg (s r : I) : infoGain s r >= 0 := by
  unfold infoGain; linarith [compose_enriches' r s]

/-- SP4. THE SILENCE PREMIUM -/
theorem silence_zero_infoGain (r : I) : infoGain (void : I) r = 0 := by
  unfold infoGain; rw [void_right']; ring
\end{lstlisting}

\section{The Chain Enrichment Theorem}

\begin{theorem}[Chains Never Lose Information]
\label{thm:chain-enrichment}
In a signaling chain $s \to m \to r$ (where $s$ is sent to intermediary $m$, who relays to receiver $r$), the receiver's final cognitive weight is at least their initial weight:
\[
\wt((r \compose m) \compose s) \geq \wt(r).
\]
\end{theorem}

\begin{proof}
We apply enrichment twice:
\begin{align*}
\wt((r \compose m) \compose s) &\geq \wt(r \compose m) && \text{(enrichment with } s\text{)} \\
&\geq \wt(r) && \text{(enrichment with } m\text{)}.
\end{align*}
Chaining these inequalities gives the result.
\end{proof}

\begin{paradox}[The Telephone Game Paradox]
In the children's game ``Telephone,'' a message is whispered from person to person, and the final message is often hilariously different from the original. The naive expectation is that information \emph{degrades} through intermediaries.

In ideatic space, the opposite is true: each intermediary adds weight. The message changes --- the emergence term at each step introduces non-additive surprise --- but the cognitive complexity of the message can only grow. The ``degradation'' in the telephone game is not a loss of information but a \emph{transformation}: emergence replaces the original structure with new structure, and the total weight monotonically increases.

This explains why rumors often become more elaborate as they spread, why conspiracy theories grow more complex over time, and why mythologies accrete detail across generations.
\end{paradox}

\section{The Double-Hearing Paradox}

\begin{theorem}[Repetition Always Enriches]
\label{thm:double-hearing}
Hearing a signal twice always produces at least as much cognitive weight as hearing it once:
\[
\wt((r \compose s) \compose s) \geq \wt(r \compose s) \geq \wt(r).
\]
\end{theorem}

\begin{proof}
Both inequalities follow directly from the enrichment axiom.
\end{proof}

\begin{interpretation}
This theorem explains why repetition is effective in persuasion, advertising, and education. Even when the content is identical, each repetition adds weight through the self-emergence $\emrg(r \compose s, s, (r \compose s) \compose s)$. The repeated signal interacts with the \emph{already-modified} receiver, producing new emergent effects.

In advertising, this is the ``mere exposure effect'': repeated exposure to a brand increases preference, even without new information. In ideatic space, this is not a psychological bias but a mathematical necessity.
\end{interpretation}

\section{Connection to Shannon's Theory}

Shannon's information theory measures information in bits --- the reduction in uncertainty about the source. A perfectly compressed message carries exactly as much information as the source entropy.

Ideatic information theory measures information in \emph{weight} --- the increase in the receiver's self-resonance. This is a fundamentally different quantity. In Shannon's theory, redundant messages carry zero additional information. In ideatic theory, redundant messages carry positive weight gain (by the Double-Hearing Paradox). This is because ideatic information is not about uncertainty reduction but about \emph{cognitive enrichment}: the interaction between new and existing ideas creates emergence, even when the ``new'' idea is a repetition.

The connection to Harsanyi's (1967--68) theory of games with incomplete information is also revealing. Harsanyi models uncertainty through type spaces, where each player has a probability distribution over the other players' types. In ideatic space, ``types'' correspond to ideas, and the type space is the ideatic space itself. But unlike Harsanyi's model, ideatic types are \emph{compositional}: two types can be composed to form a new type, and this composition is enriching. This means that learning about another player's type (by hearing their signal) always increases your cognitive weight --- a much stronger conclusion than anything in Harsanyi's framework.


% ================================================================
\chapter{Coalition Paradoxes: When Adding Members Hurts}
% ================================================================

\epigraph{In union there is strength.}{Aesop}

\section{The Naive Expectation}

In cooperative game theory, a coalition's value is determined by a \emph{characteristic function} $v: 2^N \to \mathbb{R}$ that assigns a real number to each subset of players. The key question is whether the game is \emph{superadditive}: $v(S \cup T) \geq v(S) + v(T)$ for disjoint $S, T$.

Many real-world games are \emph{not} superadditive. Adding a new member to a committee can reduce its effectiveness (too many cooks). Adding a new firm to a cartel can reduce total profit (increased coordination costs). Adding a new lane to a highway can increase travel time (Braess's paradox).

The naive expectation for ideatic space: surely some compositions reduce value? After all, conflicting ideas should cancel out.

The reality: in ideatic space, coalitions are \emph{always} superadditive in weight. Adding a member \emph{never} reduces total coalition weight.

\section{The Total Cooperation Surplus}

\begin{definition}[Total Cooperation Surplus]
The total surplus of a two-person coalition $(a, b)$ is:
\[
\Sigma(a,b) = S(a,b) + S(b,a),
\]
where $S(a,b) = \wt(a \compose b) - \wt(a)$ is the cooperation surplus from $a$'s perspective.
\end{definition}

\begin{theorem}[Total Surplus Is Non-Negative and Symmetric]
\label{thm:total-surplus}
For any ideas $a, b$:
\begin{enumerate}
\item $\Sigma(a,b) = \Sigma(b,a)$ \quad (symmetry).
\item $\Sigma(a,b) \geq 0$ \quad (non-negativity).
\item $\Sigma(a,b) = \wt(a \compose b) + \wt(b \compose a) - \wt(a) - \wt(b)$ \quad (explicit formula).
\end{enumerate}
\end{theorem}

\begin{proof}
(1) follows from the commutativity of addition: $S(a,b) + S(b,a) = S(b,a) + S(a,b)$.

(2) follows from enrichment: $S(a,b) \geq 0$ and $S(b,a) \geq 0$, so their sum is non-negative.

(3) is a direct expansion of the definitions.
\end{proof}

\begin{lstlisting}
/-- CP2. TOTAL SURPLUS IS SYMMETRIC -/
theorem totalCoopSurplus_symm (a b : I) :
    totalCoopSurplus a b = totalCoopSurplus b a := by
  unfold totalCoopSurplus; ring

/-- CP3. TOTAL SURPLUS IS NON-NEGATIVE -/
theorem totalCoopSurplus_nonneg (a b : I) :
    totalCoopSurplus a b >= 0 := by
  unfold totalCoopSurplus
  have h1 := cooperationSurplus_nonneg a b
  have h2 := cooperationSurplus_nonneg b a
  linarith
\end{lstlisting}

\section{The Enrichment Chain}

\begin{theorem}[Monotone Coalition Growth]
\label{thm:enrichment-chain}
For any three ideas $a, b, c$:
\[
\wt(a) \leq \wt(a \compose b) \leq \wt((a \compose b) \compose c).
\]
Adding members to a coalition can \emph{never} decrease total weight.
\end{theorem}

\begin{proof}
Both inequalities follow from the enrichment axiom:
\begin{itemize}
\item $\wt(a \compose b) \geq \wt(a)$: composition with $b$ enriches.
\item $\wt((a \compose b) \compose c) \geq \wt(a \compose b)$: composition with $c$ enriches.
\end{itemize}
\end{proof}

\begin{paradox}[The Impossible Tragedy of the Commons]
In Hardin's (1968) tragedy of the commons, adding users to a shared resource degrades it for everyone. Each additional cow on the commons reduces the grass available to all.

In ideatic space, there is no tragedy of the commons --- at least not in weight terms. Every addition enriches. The ``commons'' (the composed idea) grows heavier with every new participant, never lighter.

But --- and this is the subtle twist --- the \emph{distribution} of surplus can be highly asymmetric. The total pie always grows, but individual slices can shrink in relative terms. This is the ideatic version of the ``inequality paradox'': growth without equity.
\end{paradox}

\section{The Void Free Rider Theorem}

\begin{theorem}[The Perfect Free Rider]
\label{thm:void-free-rider}
When one coalition member is void:
\begin{enumerate}
\item $S(a, \void) = 0$: the non-void member gains nothing from the free rider.
\item $S(\void, a) = \wt(a)$: the free rider gains the \emph{entire} coalition value.
\end{enumerate}
\end{theorem}

\begin{proof}
(1) By right identity: $a \compose \void = a$, so $\wt(a \compose \void) = \wt(a)$, giving surplus $0$.

(2) By left identity: $\void \compose a = a$, so $\wt(\void \compose a) = \wt(a)$. The void's individual weight is $\wt(\void) = 0$ (by the void resonance axiom). Thus the surplus is $\wt(a) - 0 = \wt(a)$.
\end{proof}

\begin{interpretation}
The void idea is the perfect free rider: it contributes nothing (surplus from $a$'s perspective is zero) but captures the entire value of the coalition. This formalizes the intuition that silence can be exploitative --- by saying nothing, the silent party forces the other to ``do all the work'' of creating meaning, while benefiting from the resulting composition.

This connects to Olson's (1965) theory of collective action: in large groups, the incentive to free-ride increases because each individual's contribution is small relative to the total. In ideatic space, void is the limiting case: zero contribution, full benefit.
\end{interpretation}

\section{The Surplus Asymmetry}

\begin{theorem}[Surplus Asymmetry Is Antisymmetric]
\label{thm:surplus-asymmetry}
The surplus asymmetry $A(a,b) = S(a,b) - S(b,a)$ satisfies $A(a,b) = -A(b,a)$.
\end{theorem}

\begin{proof}
Direct computation: $A(a,b) = S(a,b) - S(b,a) = -(S(b,a) - S(a,b)) = -A(b,a)$.
\end{proof}

\begin{remark}
This theorem says that the unfairness of cooperation is a zero-sum game in itself. Whatever advantage one party gains from the order of composition, the other party loses by exactly the same amount. There is no ``win-win'' in the distribution of surplus --- only in the creation of surplus.

This connects to Nash's (1950) bargaining solution, where the bargaining problem is to divide a surplus that both parties have created. Nash showed that the unique solution satisfying certain axioms divides the surplus in proportion to the players' disagreement points. In ideatic space, the ``disagreement point'' is the individual weight $\wt(a)$, and the surplus is always non-negative. But unlike Nash's model, the surplus depends on the \emph{order} of composition, introducing an asymmetry that Nash's symmetric framework cannot capture.
\end{remark}


% ================================================================
\chapter{The Non-Commutativity of Negotiation}
% ================================================================

\epigraph{In any negotiation, the one who cares less has more power.}{Anonymous}

\section{The Naive Expectation}

In classical game theory, simultaneous-move games and sequential-move games are analyzed with different solution concepts (Nash equilibrium vs.\ subgame-perfect equilibrium), but the \emph{total value} of the game is often independent of the order of play. In many settings, the Coase theorem (1960) suggests that the initial allocation of rights doesn't matter for efficiency --- parties will negotiate to the efficient outcome regardless.

In ideatic space, the order of negotiation \emph{fundamentally} changes the outcome. Composition is not commutative, and this non-commutativity has deep strategic consequences.

\section{The Negotiation Asymmetry}

\begin{definition}[Negotiation Asymmetry]
The negotiation asymmetry between $a$ and $b$ is:
\[
N(a,b) = \wt(a \compose b) - \wt(b \compose a).
\]
\end{definition}

\begin{theorem}[Negotiation Antisymmetry]
\label{thm:negotiation-antisymm}
For all $a, b$: $N(a,b) = -N(b,a)$.
\end{theorem}

\begin{proof}
$N(a,b) = \wt(a \compose b) - \wt(b \compose a) = -(\wt(b \compose a) - \wt(a \compose b)) = -N(b,a)$.
\end{proof}

\begin{paradox}[The Perfect Zero-Sum Structure]
The negotiation asymmetry is \emph{perfectly} zero-sum: whatever advantage one party gains from going first, the other party loses by exactly the same amount. This is true for \emph{every} pair of ideas, not just those in explicit competition.

This is counter-intuitive because you'd expect some negotiations to be ``more symmetric'' than others in the sense that both parties benefit from a particular ordering. But the antisymmetry theorem says the structural advantage is always perfectly antisymmetric --- even when both orderings create value (which they always do, by enrichment).
\end{paradox}

\begin{lstlisting}
/-- DP1. NEGOTIATION ANTISYMMETRY -/
theorem negotiation_antisymm (a b : I) :
    negotiationAsym a b = -(negotiationAsym b a) := by
  unfold negotiationAsym; ring
\end{lstlisting}

\section{The Total Negotiation Value Bound}

\begin{theorem}[Negotiation Always Creates Value]
\label{thm:total-negotiation}
The total negotiation value always exceeds the sum of individual weights:
\[
\wt(a \compose b) + \wt(b \compose a) \geq \wt(a) + \wt(b).
\]
\end{theorem}

\begin{proof}
By enrichment: $\wt(a \compose b) \geq \wt(a)$ and $\wt(b \compose a) \geq \wt(b)$. Adding these two inequalities gives the result.
\end{proof}

\begin{interpretation}
Even adversarial negotiations create net value. The total weight from both possible orderings always exceeds the sum of what the parties bring individually. This means that \emph{negotiation itself} is productive, regardless of the outcome.

This connects to Myerson and Satterthwaite's (1983) impossibility theorem, which shows that efficient bilateral trade is impossible under asymmetric information. In ideatic space, ``efficient trade'' (in weight terms) is not just possible but \emph{guaranteed}: every composition enriches. The Myerson-Satterthwaite impossibility is about information rents, not about value creation. In ideatic space, value is always created; the question is how it is distributed.
\end{interpretation}

\section{Commutativity Implies Fairness}

\begin{theorem}[Commutativity Equals Fairness]
\label{thm:comm-fair}
If $a \compose b = b \compose a$, then $N(a,b) = 0$.
\end{theorem}

\begin{proof}
If $a \compose b = b \compose a$, then $\wt(a \compose b) = \wt(b \compose a)$, so $N(a,b) = 0$.
\end{proof}

\begin{corollary}
Fairness (zero negotiation asymmetry) is equivalent to commutativity of the specific composition. This means that ``fair'' negotiations are exactly those where the order of composition doesn't matter --- where the ideas combine symmetrically.
\end{corollary}

\begin{remark}
This is a startling connection between algebra and ethics. Commutativity --- a purely structural property of the monoid operation --- is equivalent to fairness --- a normative property of the negotiation outcome. In philosophical terms, this suggests that fairness is not an extrinsic value imposed on negotiations but an intrinsic property of certain algebraic structures.
\end{remark}

\section{The First-Mover Premium Decomposition}

\begin{theorem}[The First-Mover Premium]
\label{thm:first-mover}
The first-mover premium $P(a,b) = [\wt(a \compose b) - \wt(a)] - [\wt(b \compose a) - \wt(b)]$ decomposes as:
\[
P(a,b) = N(a,b) + (\wt(b) - \wt(a)).
\]
\end{theorem}

\begin{proof}
Expand:
\begin{align*}
P(a,b) &= \wt(a \compose b) - \wt(a) - \wt(b \compose a) + \wt(b) \\
&= [\wt(a \compose b) - \wt(b \compose a)] + [\wt(b) - \wt(a)] \\
&= N(a,b) + (\wt(b) - \wt(a)).
\end{align*}
\end{proof}

\begin{paradox}[The Paradox of the Weak First Mover]
The first-mover premium has \emph{two} components: the negotiation asymmetry $N(a,b)$ and the weight differential $\wt(b) - \wt(a)$. Even when the negotiation asymmetry is zero (fair composition), the first-mover premium can be positive if the second mover is heavier than the first.

This means that in a ``fair'' negotiation (commutative composition), the \emph{lighter} party benefits more from going first. The weak benefit from the first-mover advantage more than the strong do.

This inverts the usual intuition. In most negotiation settings, the stronger party benefits from going first (they set the anchor, frame the discussion). In ideatic space, when composition commutes, the weaker party gains more from the structural advantage of going first, because the weight differential contributes positively to their premium.
\end{paradox}

\begin{lstlisting}
/-- DP9. THE FIRST-MOVER PREMIUM DECOMPOSITION -/
theorem fmPremium_decomposition (a b : I) :
    fmPremium a b = negotiationAsym a b + (rs b b - rs a a) := by
  unfold fmPremium weightSurplus negotiationAsym; ring
\end{lstlisting}


% ================================================================
\chapter{Emergence in Zero-Sum Games: Something From Nothing}
% ================================================================

\epigraph{The whole is greater than the sum of its parts.}{Aristotle (attributed)}

\section{The Naive Expectation}

Zero-sum games are the paradigmatic case of pure competition: one player's gain is another's loss. Von Neumann's minimax theorem (1928) guarantees the existence of optimal strategies in such games, and the total payoff is always zero.

In ideatic space, something extraordinary happens: even when the \emph{resonance} between two ideas is zero (they are orthogonal), their composition still creates positive weight. \emph{Something emerges from nothing.}

\section{The Zero-Sum Emergence Paradox}

\begin{theorem}[Commutative Emergence Symmetry]
\label{thm:zero-sum-emergence}
If $a \compose b = b \compose a$ (composition commutes), then:
\[
\emrg(a,b,c) = \emrg(b,a,c) \quad \text{for all } c.
\]
Emergence is symmetric when composition commutes. But it need NOT be zero.
\end{theorem}

\begin{proof}
By definition, $\emrg(a,b,c) = \operatorname{rs}(a \compose b, c) - \operatorname{rs}(a,c) - \operatorname{rs}(b,c)$. If $a \compose b = b \compose a$, then
\[
\emrg(a,b,c) = \operatorname{rs}(a \compose b, c) - \operatorname{rs}(a,c) - \operatorname{rs}(b,c) = \operatorname{rs}(b \compose a, c) - \operatorname{rs}(a,c) - \operatorname{rs}(b,c) = \emrg(b,a,c).
\]
But this says nothing about whether $\emrg(a,b,c)$ is zero --- only that it equals $\emrg(b,a,c)$.
\end{proof}

\begin{paradox}[Something From Nothing]
When $a \compose b = b \compose a$, the composition is ``symmetric'' --- neither party has an ordering advantage. You might expect this to mean the composition is ``boring'' --- purely additive, with no emergence.

But emergence can be nonzero even for commutative compositions! The composition $a \compose b$ can resonate with observer $c$ in ways that $a$ and $b$ individually do not, even when the order doesn't matter.

This is the ideatic analog of constructive interference in physics: two waves can combine to produce an amplitude greater than either alone, even when neither wave ``knows'' about the other.
\end{paradox}

\section{Self-Composition Always Enriches}

\begin{theorem}[Echo Chamber Effect]
For any idea $a$:
\[
\wt(a \compose a) \geq \wt(a).
\]
\end{theorem}

\begin{proof}
Direct application of the enrichment axiom with $b = a$.
\end{proof}

\begin{interpretation}
Repeating an idea to yourself always increases your cognitive weight. The ``echo chamber effect'' is not a social phenomenon but a mathematical theorem. Self-reinforcement is algebraically guaranteed.

This connects to Sunstein's (2001) work on group polarization: when like-minded individuals discuss an issue, they tend to move toward a more extreme position. In ideatic space, this is a consequence of self-composition enrichment: composing similar ideas always increases weight, driving the group toward heavier (more extreme) positions.
\end{interpretation}

\section{The Attribution Conservation Law}

\begin{theorem}[Attribution Is Inflationary]
\label{thm:attribution-conservation}
The sum of marginal contributions in both orderings equals:
\[
M(a,b,c) + M(b,a,c) = \operatorname{rs}(a \compose b, c) + \emrg(b,a,c),
\]
where $M(a,b,c) = \operatorname{rs}(a \compose b, c) - \operatorname{rs}(b, c)$ is $a$'s marginal contribution to the pair $(a,b)$ as seen by $c$.
\end{theorem}

\begin{proof}
Expand both marginal contributions:
\begin{align*}
M(a,b,c) + M(b,a,c) &= [\operatorname{rs}(a \compose b, c) - \operatorname{rs}(b,c)] + [\operatorname{rs}(b \compose a, c) - \operatorname{rs}(a,c)] \\
&= \operatorname{rs}(a \compose b, c) + \operatorname{rs}(b \compose a, c) - \operatorname{rs}(a,c) - \operatorname{rs}(b,c) \\
&= \operatorname{rs}(a \compose b, c) + [\operatorname{rs}(b \compose a, c) - \operatorname{rs}(b,c) - \operatorname{rs}(a,c)] \\
&= \operatorname{rs}(a \compose b, c) + \emrg(b,a,c).
\end{align*}
\end{proof}

\begin{paradox}[The Impossibility of Fair Attribution]
The sum of marginal contributions does NOT equal the total value $\operatorname{rs}(a \compose b, c)$ --- it exceeds it by $\emrg(b,a,c)$. This means that if you try to assign credit based on marginal contributions (as Shapley values do), you will \emph{over-allocate}: the sum of credits exceeds the total value.

This is the ideatic impossibility of fair attribution. In cooperative game theory, the Shapley value provides a unique ``fair'' allocation under certain axioms. In ideatic space, no such allocation exists whenever emergence is nonzero --- the credit pie is always bigger than the value pie.

The over-allocation equals exactly the emergence in the reverse order: $\emrg(b,a,c)$. This is the ``double-counting'' of credit that arises because both $a$ and $b$ claim credit for the emergent surplus, but emergence belongs to neither.
\end{paradox}


% --- Part II ---
\part{The Paradoxes of Cooperation and Competition}

% ================================================================
% Volume II-B, Chapters 6--10
% The Paradoxes of Cooperation and Competition
% ================================================================

% ================================================================
\chapter{The Braess Paradox of Ideas: When More Knowledge Hurts}
% ================================================================

\epigraph{Knowledge is power. Knowledge shared is power multiplied.}{Robert Noyce}

\section{The Naive Expectation}

Braess's paradox (1968) is the classic result that adding a road to a traffic network can increase total travel time. It arises because selfish routing leads to a worse equilibrium when more options are available.

For ideas, the naive expectation is similar: surely, giving someone \emph{more} information should always help, or at worst be neutral? After all, a rational agent can always ignore new information.

In ideatic space, you \emph{cannot} ignore new information. Composition is irreversible. And this leads to a striking analog of Braess's paradox.

\section{The Order Information Asymmetry}

\begin{theorem}[The Framing Effect Theorem]
\label{thm:framing-effect}
The difference between hearing $s_1$ then $s_2$ versus $s_2$ then $s_1$ is:
\begin{align*}
\operatorname{rs}((r \compose s_1) &\compose s_2, c) - \operatorname{rs}((r \compose s_2) \compose s_1, c) \\
&= [\emrg(r, s_1, c) + \emrg(r \compose s_1, s_2, c)] \\
&\quad - [\emrg(r, s_2, c) + \emrg(r \compose s_2, s_1, c)].
\end{align*}
\end{theorem}

\begin{proof}
Apply the decomposition identity $\operatorname{rs}(a \compose b, c) = \operatorname{rs}(a,c) + \operatorname{rs}(b,c) + \emrg(a,b,c)$ four times:

For $\operatorname{rs}((r \compose s_1) \compose s_2, c)$:
\begin{align*}
&= \operatorname{rs}(r \compose s_1, c) + \operatorname{rs}(s_2, c) + \emrg(r \compose s_1, s_2, c) \\
&= [\operatorname{rs}(r, c) + \operatorname{rs}(s_1, c) + \emrg(r, s_1, c)] + \operatorname{rs}(s_2, c) + \emrg(r \compose s_1, s_2, c).
\end{align*}

Similarly for $\operatorname{rs}((r \compose s_2) \compose s_1, c)$:
\begin{align*}
&= \operatorname{rs}(r, c) + \operatorname{rs}(s_2, c) + \emrg(r, s_2, c) + \operatorname{rs}(s_1, c) + \emrg(r \compose s_2, s_1, c).
\end{align*}

Subtracting, the $\operatorname{rs}(r,c)$, $\operatorname{rs}(s_1,c)$, and $\operatorname{rs}(s_2,c)$ terms cancel, leaving the emergence chain differential.
\end{proof}

\begin{lstlisting}
/-- SP9. THE FRAMING EFFECT THEOREM -/
theorem framing_effect_theorem (s1 s2 r c : I) :
    rs (compose (compose r s1) s2) c -
    rs (compose (compose r s2) s1) c =
    (emergence r s1 c + emergence (compose r s1) s2 c) -
    (emergence r s2 c + emergence (compose r s2) s1 c) := by
  have h1 := rs_compose_eq r s1 c
  have h2 := rs_compose_eq (compose r s1) s2 c
  have h3 := rs_compose_eq r s2 c
  have h4 := rs_compose_eq (compose r s2) s1 c
  linarith
\end{lstlisting}

\begin{paradox}[The Braess Paradox of Ideas]
While both orderings of two signals \emph{both} enrich the receiver (by the Dual Order Enrichment theorem), they may enrich to \emph{different} levels. One ordering might produce significantly more weight than the other.

This means that the \emph{order} in which you learn things matters as much as \emph{what} you learn. In educational settings, presenting concepts in the ``wrong'' order can produce a less enriched cognitive state than the optimal order, even though both orders present the same material.

This is Braess's paradox for ideas: the ``route'' through information space (which idea you encounter first) affects the final destination, and some routes are strictly worse than others --- even though all routes increase total weight.
\end{paradox}

\section{The Dual Order Enrichment}

\begin{theorem}[Both Orders Enrich]
\label{thm:dual-enrichment}
For any signals $s_1, s_2$ and receiver $r$:
\begin{enumerate}
\item $\wt((r \compose s_1) \compose s_2) \geq \wt(r)$.
\item $\wt((r \compose s_2) \compose s_1) \geq \wt(r)$.
\end{enumerate}
Both orderings enrich the receiver, even though they may enrich to different levels.
\end{theorem}

\begin{proof}
For (1): $\wt((r \compose s_1) \compose s_2) \geq \wt(r \compose s_1) \geq \wt(r)$ by two applications of enrichment.

(2) is symmetric: $\wt((r \compose s_2) \compose s_1) \geq \wt(r \compose s_2) \geq \wt(r)$.
\end{proof}

\begin{interpretation}
This is the ``silver lining'' of the Braess paradox: while some orderings are better than others, NO ordering is worse than not hearing the information at all. Every path through information space leads to enrichment. You can choose suboptimally, but you cannot choose destructively.

This connects to the ``no regret'' framework in online learning (Hannan, 1957; Blackwell, 1956): an algorithm that ensures bounded regret can learn from any sequence of information without being ``hurt'' by adversarial ordering. In ideatic space, regret is bounded by construction --- the enrichment axiom guarantees it.
\end{interpretation}

\section{Connection to Selten's Subgame Perfection}

Selten's (1965) concept of subgame-perfect equilibrium refines Nash equilibrium for sequential games by requiring that strategies be optimal at every point in the game tree, not just at the beginning. The key insight is that sequential structure matters: a player who moves later has different information than one who moves earlier.

In ideatic space, the sequential structure is captured by the composition order. The Framing Effect Theorem shows that the emergence chain --- the sequence of creative surpluses at each step --- determines the final outcome. Selten's subgame perfection corresponds to choosing the composition order that maximizes one's emergence at each step. But unlike classical subgame perfection, both orderings are always enriching --- there is no ``threat'' that is not credible, because every composition enriches.


% ================================================================
\chapter{Fixed Points of Strategic Composition}
% ================================================================

\epigraph{Give me a place to stand, and I shall move the Earth.}{Archimedes}

\section{The Naive Expectation}

Fixed-point theorems are among the most powerful tools in game theory. Brouwer's fixed-point theorem guarantees the existence of Nash equilibria in finite games. Kakutani's fixed-point theorem extends this to mixed strategies. Tarski's fixed-point theorem applies to lattices.

The naive expectation: ideatic spaces should have fixed points of composition, and these should correspond to equilibria. But the actual fixed-point structure is richer and stranger than expected.

\section{Void as the Universal Fixed Point}

\begin{theorem}[Void Is a Fixed Point]
\label{thm:void-fixed-point}
The void idea is a fixed point of composition with anything:
\[
\void \compose a = a = a \compose \void \quad \text{for all } a.
\]
\end{theorem}

\begin{proof}
This is the left and right identity axioms.
\end{proof}

\begin{interpretation}
Void is not just ``silence'' --- it is the \emph{identity element} of the compositional monoid. In game-theoretic terms, ``doing nothing'' is always a fixed point. This connects to the result from Volume~II that void is always a Nash equilibrium: $(void, void)$ is an equilibrium because silence is always a best response to silence.

But void is a degenerate fixed point --- it represents the absence of interaction. The interesting question is whether there are \emph{non-trivial} fixed points.
\end{interpretation}

\section{The Weight Ratchet and Convergence}

\begin{theorem}[The Weight Ratchet]
\label{thm:weight-ratchet}
For any idea $a$, the sequence $\wt(a), \wt(a \compose a), \wt((a \compose a) \compose a), \ldots$ is non-decreasing:
\[
\wt(\compN{a}{n+1}) \geq \wt(\compN{a}{n}) \quad \text{for all } n.
\]
\end{theorem}

\begin{proof}
$\compN{a}{n+1} = \compN{a}{n} \compose a$, so $\wt(\compN{a}{n+1}) \geq \wt(\compN{a}{n})$ by enrichment.
\end{proof}

\begin{paradox}[The Ratchet Cannot Be Reversed]
Once you compose with an idea, your weight can never return to its previous level through further composition. There is no ``inverse'' of an idea in the weight sense. The ratchet only goes one way.

This means that in any strategic interaction modeled as iterated composition, \emph{commitment is irreversible}. A player who has composed with $a$ cannot ``undo'' this composition --- they can only compose further. This formalizes the game-theoretic concept of sunk costs: investments in ideas are permanent.
\end{paradox}

\section{Step Emergence and Telescoping}

\begin{theorem}[Weight Telescoping]
\label{thm:telescoping}
The weight of $\compN{a}{n+1}$ equals the sum of step emergences:
\[
\wt(\compN{a}{n+1}) = \sum_{i=0}^{n} \left[\wt(\compN{a}{i+1}) - \wt(\compN{a}{i})\right],
\]
where each step emergence $\Delta_i = \wt(\compN{a}{i+1}) - \wt(\compN{a}{i}) \geq 0$.
\end{theorem}

\begin{proof}
This is a telescoping sum: the intermediate terms cancel, and we use $\wt(\compN{a}{0}) = \wt(\void) = 0$.
\end{proof}

\begin{interpretation}
Weight is \emph{cumulative}: it is the sum of all the ``surprises'' (step emergences) from each iteration. This means that the total cognitive weight of an idea after $n$ repetitions is exactly the accumulated creativity from each step.

In economic terms, this is like compound interest, except the ``interest rate'' (step emergence) can vary from step to step and is always non-negative. The weight grows, but the growth rate is not constant.
\end{interpretation}


% ================================================================
\chapter{The Information Bound: Fundamental Limits on Signaling}
% ================================================================

\epigraph{Information is the resolution of uncertainty.}{Claude Shannon}

\section{The Naive Expectation}

Shannon's channel capacity theorem says that every communication channel has a maximum rate at which information can be reliably transmitted. Beyond this rate, errors become inevitable.

In ideatic space, there is an analogous fundamental bound: the emergence bound, which limits how much ``surprise'' a signal can produce.

\section{The Fundamental Manipulation Bound}

\begin{theorem}[The Manipulation Bound]
\label{thm:manipulation-bound}
For any signal $s$ and receiver $r$:
\[
\emrg(r, s, r)^2 \leq \wt(r \compose s) \cdot \wt(r).
\]
The squared ``surprise'' that the receiver experiences (measured by emergence with themselves as observer) is bounded by the product of the composed weight and the receiver's own weight.
\end{theorem}

\begin{proof}
This is a direct application of the emergence bound axiom with $a = r$, $b = s$, $c = r$.
\end{proof}

\begin{lstlisting}
/-- SP10. THE FUNDAMENTAL MANIPULATION BOUND -/
theorem manipulation_cost_bound (s r : I) :
    (emergence r s r)^2 <=
    rs (compose r s) (compose r s) * rs r r :=
  emergence_bounded' r s r
\end{lstlisting}

\begin{paradox}[Manipulation Is Expensive]
The manipulation bound says that to produce a large emergence (a big ``surprise'' in the receiver), you need \emph{heavy} ideas. Specifically:
\begin{itemize}
\item If the signal is weak ($\wt(r \compose s)$ is small), the surprise must be small.
\item If the receiver is weak ($\wt(r)$ is small), the surprise must be small.
\item Large surprises require BOTH a heavy composed result AND a heavy receiver.
\end{itemize}

This explains why propaganda targets the cognitively ``heavy'' (those with strong existing beliefs): lightweight minds produce only small emergence, so the manipulation effect is small. The irony is that the most manipulable people are those with the strongest existing beliefs --- they have the most weight to ``resonate with.''
\end{paradox}

\section{The Polarization Bound}

\begin{theorem}[Emergence Is Bounded By Weights]
\label{thm:polarization-bound}
For any ideas $a$, $b$ and observer $c$:
\[
\emrg(a,b,c)^2 \leq \wt(a \compose b) \cdot \wt(c).
\]
\end{theorem}

\begin{proof}
This is the emergence bound axiom.
\end{proof}

\begin{interpretation}
This theorem places a fundamental ceiling on polarization. The emergence between two ideas, as observed by a third party, cannot exceed the geometric mean of the composed weight and the observer's weight. This means:

\begin{enumerate}
\item \textbf{Light observers see small effects}: If $\wt(c)$ is small (the observer has few existing beliefs), emergence is bounded to be small. Newcomers to a debate see less polarization than veterans.

\item \textbf{Light compositions produce small effects}: If $\wt(a \compose b)$ is small (the combined idea is lightweight), emergence is bounded. Trivial debates produce little emergence.

\item \textbf{The ceiling is multiplicative}: The bound is the \emph{product} $\wt(a \compose b) \cdot \wt(c)$, not the sum. This means that both the composition and the observer must be heavy for large emergence. A heavy composition observed by a lightweight observer (or vice versa) produces bounded emergence.
\end{enumerate}

This connects to Aumann's (1976) ``agreeing to disagree'' theorem: rational agents with common priors cannot agree to disagree. In ideatic space, ``disagreement'' (negative emergence) is bounded by the polarization bound. The maximum possible disagreement is limited by the agents' weights.
\end{interpretation}

\section{The No-Diminishing-Returns Theorem}

\begin{theorem}[Every Signal Adds Value]
\label{thm:no-diminishing-returns}
The second-order information gain is always non-negative:
\[
\wt((r \compose s_1) \compose s_2) - \wt(r \compose s_1) \geq 0.
\]
\end{theorem}

\begin{proof}
Direct application of enrichment to the pair $((r \compose s_1), s_2)$.
\end{proof}

\begin{paradox}[No Diminishing Returns to Information]
In economics, diminishing marginal returns is one of the most fundamental principles: the tenth unit of a good provides less utility than the first. You'd expect the same for information: the tenth fact about a topic should be less valuable than the first.

In ideatic space, there are NO diminishing returns to information (in weight terms). The $n$th signal always provides non-negative weight gain, regardless of how many signals preceded it. This is because each new signal interacts with the \emph{already-enriched} receiver, producing new emergence that is independent of previous emergences.

This explains information addiction: each new piece of information is independently rewarding (in cognitive weight terms), creating a positive feedback loop with no natural satiation point.
\end{paradox}


% ================================================================
\chapter{Evolutionary Paradoxes: Why the Best Strategy Loses}
% ================================================================

\epigraph{It is not the strongest of the species that survives, nor the most intelligent that survives. It is the one that is most adaptable to change.}{Charles Darwin (attributed)}

\section{The Naive Expectation}

In evolutionary game theory (Maynard Smith and Price, 1973), the fittest strategy --- the one with the highest payoff --- tends to dominate the population. The replicator dynamics ensures that strategies with above-average fitness increase in frequency.

In ideatic space, evolution works through \emph{composition} rather than replication. Ideas don't copy themselves; they compose with their context, producing offspring that differ from the parent via emergence. This leads to several counter-intuitive results.

\section{The Complexity Ratchet}

\begin{theorem}[Cultural Complexity Never Decreases]
\label{thm:complexity-ratchet}
For any idea $a$ and context $b$, the weight of the drifted idea after $n$ generations is non-decreasing:
\[
\wt(\text{drift}(a,b,n+1)) \geq \wt(\text{drift}(a,b,n)),
\]
where $\text{drift}(a,b,0) = a$ and $\text{drift}(a,b,n+1) = \text{drift}(a,b,n) \compose b$.
\end{theorem}

\begin{proof}
$\text{drift}(a,b,n+1) = \text{drift}(a,b,n) \compose b$, so by enrichment:
\[
\wt(\text{drift}(a,b,n+1)) = \wt(\text{drift}(a,b,n) \compose b) \geq \wt(\text{drift}(a,b,n)).
\]
\end{proof}

\begin{lstlisting}
/-- EP1. THE COMPLEXITY RATCHET -/
theorem complexity_ratchet (a b : I) (n m : N) (h : n <= m) :
    rs (drift a b m) (drift a b m) >=
    rs (drift a b n) (drift a b n) :=
  drift_nondecreasing a b n m h
\end{lstlisting}

\begin{paradox}[The Impossibility of Simplification]
In biological evolution, organisms can become simpler: parasites lose organs, cave fish lose eyes, endosymbionts lose genes. Simplification is a real evolutionary strategy.

In ideatic evolution, simplification is \emph{mathematically impossible}. Every generation of cultural drift adds weight; no generation can remove it. Cultural complexity can only grow.

This explains the observation that cultural artifacts (languages, legal systems, bureaucracies, software) tend to grow more complex over time but rarely simplify. It's not just organizational inertia or path dependence --- it's a mathematical theorem. The enrichment axiom guarantees that cultural complexity is a one-way ratchet.
\end{paradox}

\section{The Extinction Impossibility}

\begin{theorem}[Ideas Cannot Go Extinct]
\label{thm:no-extinction}
If $a \neq \void$, then $\text{drift}(a, b, n) \neq \void$ for all $n$ and $b$.
\end{theorem}

\begin{proof}
By induction on $n$.

\emph{Base case}: $\text{drift}(a,b,0) = a \neq \void$ by hypothesis.

\emph{Inductive step}: Assume $\text{drift}(a,b,k) \neq \void$. Then $\text{drift}(a,b,k+1) = \text{drift}(a,b,k) \compose b$. Since $\text{drift}(a,b,k) \neq \void$, we have $\text{drift}(a,b,k) \compose b \neq \void$ by the non-degeneracy of composition: if the left factor is non-void, the composition is non-void.
\end{proof}

\begin{lstlisting}
/-- EP5. THE EXTINCTION IMPOSSIBILITY -/
theorem no_extinction (a b : I) (h : a != void) (n : N) :
    drift a b n != void := by
  induction n with
  | zero => simp [drift]; exact h
  | succ k ih =>
    simp [drift]
    exact compose_ne_void_of_left (drift a b k) b ih
\end{lstlisting}

\begin{paradox}[Immortal Ideas]
In biology, species go extinct. Most species that have ever lived are now extinct. Extinction is the norm, survival the exception.

In ideatic space, extinction is \emph{impossible}. Once an idea is instantiated (is non-void), no amount of cultural drift can reduce it to void. The idea may change beyond recognition through emergence, but it can never become ``nothing.''

This has profound implications for cultural evolution: ideas are immortal in a way that organisms are not. A cultural meme, once created, persists indefinitely in some form. It may be transformed, reinterpreted, or buried under layers of composition, but it can never be truly extinguished.

This connects to Dawkins's (1976) concept of the ``meme'' but goes further: while Dawkins proposed memes as cultural replicators analogous to genes, ideatic theory shows that memes are \emph{stronger} than genes --- they cannot go extinct through any natural process of cultural transmission.
\end{paradox}

\section{Mutation Always Enriches}

\begin{theorem}[All Mutations Are Beneficial]
\label{thm:beneficial-mutations}
The offspring of idea $a$ in context $b$ always has at least as much weight as the parent:
\[
\wt(a \compose b) \geq \wt(a).
\]
\end{theorem}

\begin{proof}
Direct application of the enrichment axiom.
\end{proof}

\begin{paradox}[No Harmful Mutations]
In biology, most mutations are neutral or harmful; beneficial mutations are rare. This is a cornerstone of population genetics.

In ideatic evolution, ALL ``mutations'' (offspring that differ from parent via emergence) are weakly beneficial. Every offspring is at least as heavy as its parent.

This inverts the central assumption of evolutionary biology. The reason is that ideatic ``mutation'' is not random perturbation of a genome --- it is structured composition with context. The emergence term introduces non-random variation that is constrained by the enrichment axiom to be non-negative.
\end{paradox}

\section{Connection to Evolutionary Game Theory}

Maynard Smith's (1982) concept of the Evolutionarily Stable Strategy (ESS) requires that the incumbent strategy be able to resist invasion by any mutant. In ideatic space, ``invasion'' corresponds to composition with a new idea, and ``resistance'' corresponds to weight maintenance.

Since composition always enriches, the incumbent idea's weight can only grow when ``invaded'' by a new idea. This means that in ideatic space, \emph{every} idea is ``resistant to invasion'' in the weight sense --- its weight can never decrease. But the relative position can change: a lighter idea composed with a heavier context can grow faster than a heavier idea composed with a lighter context.

This connects to the ``Red Queen'' hypothesis of Van Valen (1973): organisms must constantly evolve just to maintain their relative position. In ideatic space, the Red Queen is not a metaphor but a theorem: even though your absolute weight can never decrease, your relative fitness changes as other ideas compose and grow.


% ================================================================
\chapter{The Price of Stability vs.\ the Price of Anarchy}
% ================================================================

\epigraph{The price of liberty is eternal vigilance.}{Thomas Jefferson (attributed)}

\section{The Naive Expectation}

In mechanism design, the ``price of anarchy'' (Koutsoupias and Papadimitriou, 1999) measures how much worse the worst Nash equilibrium is compared to the social optimum. The ``price of stability'' (Anshelevich et al., 2008) measures how much worse the \emph{best} Nash equilibrium is.

In classical games, these can be arbitrarily large. In ideatic space, the situation is remarkably different.

\section{The Void Equilibrium and Its Value}

From Volume~II, we know that $(\void, \void)$ is always a Nash equilibrium. This equilibrium has zero social welfare: both players contribute nothing and receive nothing.

Any non-void equilibrium $(s, r)$ has communication surplus $\operatorname{cs}(s,r) \geq 0$ (by the Nash surplus theorem from Volume~II). This means that the ``price of stability'' is zero: the best equilibrium is always at least as good as the void equilibrium.

\begin{theorem}[Zero Price of Stability]
In any meaning game, the best Nash equilibrium has non-negative surplus.
\end{theorem}

\begin{proof}
The void equilibrium has surplus zero. Any other equilibrium has non-negative surplus by the Nash surplus theorem. Therefore the best equilibrium has surplus $\geq 0$.
\end{proof}

\section{The Network Weight Ratchet}

\begin{theorem}[Network Weight Never Decreases]
For any network node $a$ and iteration count $n$:
\[
\wt(\compN{a}{n+1}) \geq \wt(\compN{a}{n}).
\]
Network weight is monotonically non-decreasing through interaction.
\end{theorem}

\begin{proof}
$\compN{a}{n+1} = \compN{a}{n} \compose a$, and enrichment gives $\wt(\compN{a}{n} \compose a) \geq \wt(\compN{a}{n})$.
\end{proof}

\begin{interpretation}
This means that in a network of interacting ideas, the total weight can only grow. There is no ``degradation'' through interaction --- each round of interaction enriches every participating node.

This connects to the ``price of anarchy'' in a novel way: if we define the social optimum as the maximum total weight achievable through composition, and the ``anarchic'' outcome as the weight resulting from uncoordinated interaction, then the price of anarchy is bounded above by 1 --- because uncoordinated interaction always enriches. The ``worst'' outcome of anarchy is still better than the initial state.
\end{interpretation}

\section{The Consensus Enrichment Paradox}

\begin{theorem}[Adding Members Always Enriches Consensus]
For any idea $a$ and any existing consensus list:
\[
\wt(\text{consensus}[a, b]) \geq \wt(\text{consensus}[a]) = \wt(a).
\]
\end{theorem}

\begin{proof}
$\text{consensus}[a, b] = a \compose b$, and enrichment gives $\wt(a \compose b) \geq \wt(a)$.
\end{proof}

\begin{paradox}[Dissent Cannot Weaken Consensus]
In social choice theory, adding a dissenting voter can change the outcome of an election (Condorcet's paradox). In organizational theory, adding a member to a committee can reduce its effectiveness.

In ideatic space, adding a member to a consensus \emph{always} enriches it. Even a member who ``opposes'' the existing consensus (has negative resonance with it) adds weight through the enrichment axiom.

This seems to contradict everyday experience. The resolution is that ``weight'' and ``agreement'' are different things. The consensus grows heavier (more complex, more informationally rich) even when it becomes less coherent. The ``weight'' includes the cognitive burden of managing disagreement, which is itself a form of enrichment.
\end{paradox}

\section{Bridge Nodes Are Strengthened}

\begin{theorem}[The Bridge Strengthening Theorem]
A bridge node connecting two clusters gains weight from both:
\[
\wt((\text{bridge} \compose \text{left}) \compose \text{right}) \geq \wt(\text{bridge}).
\]
\end{theorem}

\begin{proof}
Two applications of enrichment: first composing with ``left'' enriches, then composing with ``right'' enriches further.
\end{proof}

\begin{interpretation}
In network theory, bridges are typically viewed as structural vulnerabilities --- if a bridge fails, the network fragments. In ideatic space, bridges are the \emph{strongest} nodes: they gain weight from both clusters they connect.

This explains why interdisciplinary researchers (bridges between academic fields) often develop the richest cognitive frameworks. Their ideas gain weight from multiple communities, making them heavier (more complex and nuanced) than specialists who interact with only one community.
\end{interpretation}


% --- Part III ---
\part{The Deep Structure of Strategic Interaction}

% ================================================================
% Volume II-B, Chapters 11--15
% The Deep Structure of Strategic Interaction
% ================================================================

% ================================================================
\chapter{Trust Cascades: Small Violations, Large Consequences}
% ================================================================

\epigraph{Trust takes years to build, seconds to break, and forever to repair.}{Anonymous}

\section{The Naive Expectation}

In repeated game theory (Fudenberg and Maskin, 1986), the folk theorem says that almost any individually rational payoff can be sustained as a Nash equilibrium in an infinitely repeated game, provided players are sufficiently patient. Trust is modeled as a discount factor: higher patience sustains more cooperative outcomes.

The naive expectation is that small trust violations should have small consequences --- proportional penalties for proportional defections. In ideatic space, this expectation is both confirmed and subverted.

\section{Trust Level Equals Cooperation Surplus}

\begin{theorem}[Trust Is Surplus]
\label{thm:trust-surplus}
The trust level of $a$ toward $b$ equals the cooperation surplus:
\[
T(a,b) = \wt(a \compose b) - \wt(a) = S(a,b) \geq 0.
\]
\end{theorem}

\begin{proof}
Both quantities equal $\wt(a \compose b) - \wt(a)$, which is non-negative by enrichment.
\end{proof}

\begin{interpretation}
Trust is not an abstract psychological quantity --- it is the \emph{cooperation surplus}. When you trust someone, what you gain is exactly the weight increase from composing with their idea. Trust and cooperation are algebraically identical.

This connects to Axelrod's (1984) ``evolution of cooperation'': in the iterated Prisoner's Dilemma, cooperation emerges through reciprocity. In ideatic space, cooperation (measured by surplus) IS trust (measured by weight gain). The two concepts collapse into one.
\end{interpretation}

\section{Betrayal Cost Is Exactly Trust}

\begin{theorem}[Betrayal Cost]
\label{thm:betrayal-cost}
The cost of betrayal (the loss from the partner defecting to void) equals the trust level:
\[
B(a,b) = \wt(a \compose b) - \wt(a \compose \void) = T(a,b).
\]
\end{theorem}

\begin{proof}
$\wt(a \compose \void) = \wt(a)$ by right identity. So $B(a,b) = \wt(a \compose b) - \wt(a) = T(a,b)$.
\end{proof}

\begin{paradox}[The Vulnerability Paradox]
Your vulnerability to betrayal is exactly your benefit from cooperation. The more you benefit from a partnership, the more you can be hurt by its dissolution. Trust and vulnerability are \emph{the same quantity}.

This formalizes the folk wisdom that ``the bigger they are, the harder they fall'': coalitions that create the most surplus are the most vulnerable to defection. There is no way to capture the benefits of cooperation without being equally exposed to the costs of betrayal.

This contradicts the popular self-help advice to ``trust but verify'': verification reduces vulnerability, but it \emph{equally} reduces the cooperation surplus. You cannot make yourself less vulnerable without also making yourself less cooperative.
\end{paradox}

\section{The Trust Asymmetry}

\begin{theorem}[Trust Asymmetry Equals Cooperation Asymmetry]
\label{thm:trust-asymmetry}
$T(a,b) - T(b,a) = A(a,b)$ (the cooperation asymmetry).
\end{theorem}

\begin{proof}
Direct expansion of definitions:
\[
T(a,b) - T(b,a) = [\wt(a \compose b) - \wt(a)] - [\wt(b \compose a) - \wt(b)] = A(a,b).
\]
\end{proof}

\begin{interpretation}
When trust is asymmetric --- one party benefits more from the partnership than the other --- the asymmetry equals the cooperation asymmetry, which equals the difference in coalition values minus the difference in individual values.

This means that ``unequal trust'' in a relationship is detectable: it manifests as different cooperation surpluses. The party who trusts more benefits more from the partnership but is also more vulnerable to betrayal.
\end{interpretation}

\section{The Cascade Irreversibility}

\begin{theorem}[Cascade Irreversibility]
An information cascade can only grow heavier:
\[
\wt((a \compose b) \compose c) \geq \wt(a \compose b) \geq \wt(a).
\]
\end{theorem}

\begin{proof}
Two applications of enrichment.
\end{proof}

\begin{interpretation}
Once a trust cascade starts (a chain of compositions), it cannot be reversed. Each link in the chain adds weight. This explains why trust collapses are so devastating: when a cascade runs in reverse (a series of betrayals), each betrayal doesn't reduce weight --- it adds the weight of the betrayal experience itself. The resulting ``distrust cascade'' is \emph{heavier} than the original trust cascade, not lighter.

This connects to Geanakoplos's (1992) work on common knowledge: trust cascades are a form of common knowledge --- each party knows that each party knows that they cooperate. In ideatic space, common knowledge is compositional weight: each layer of ``I know that you know'' adds weight through composition.
\end{interpretation}


% ================================================================
\chapter{The Manipulation Bound: How Much Can You Distort?}
% ================================================================

\epigraph{You can fool all the people some of the time, and some of the people all the time, but you cannot fool all the people all the time.}{Abraham Lincoln (attributed)}

\section{The Naive Expectation}

In mechanism design (Myerson, 1981), the revelation principle says that any equilibrium outcome of any mechanism can be replicated by a truthful mechanism. This suggests that manipulation has no fundamental power: anything achievable through manipulation is achievable through honesty.

But this assumes that the mechanism designer has full control over the information structure. In ideatic space, manipulation is constrained by a fundamental \emph{physical} bound, not just an incentive-compatibility constraint.

\section{The Emergence Bound as Manipulation Limit}

\begin{theorem}[The Fundamental Manipulation Bound]
\label{thm:manipulation-bound-2}
For any signal $s$ and receiver $r$:
\[
|\emrg(r, s, r)| \leq \sqrt{\wt(r \compose s) \cdot \wt(r)}.
\]
The ``surprise'' from any signal is bounded by the geometric mean of the composed weight and the receiver's weight.
\end{theorem}

\begin{proof}
The emergence bound axiom gives $\emrg(r,s,r)^2 \leq \wt(r \compose s) \cdot \wt(r)$. Taking square roots gives the result.
\end{proof}

\begin{paradox}[Lightweight Lies Are Impotent]
To produce a large surprise (a large emergence), you need HEAVY ideas --- both the composed result and the receiver must be heavy. A lightweight lie interacting with a lightweight receiver produces only a small effect.

This inverts the common intuition that lies are ``easy'' and truth is ``hard.'' In ideatic space, EFFECTIVE lies are expensive: they require heavy ideas. Cheap lies produce only cheap effects.

Specifically:
\begin{enumerate}
\item If $\wt(r) \approx 0$ (the receiver has few beliefs), manipulation is nearly impossible. You cannot manipulate someone who believes nothing.
\item If $\wt(r \compose s)$ is small (the lie is lightweight), the manipulation effect is small. Elaborate lies are necessary for large effects.
\item The bound is \emph{multiplicative}: both factors must be large. A heavy lie targeting a light receiver (or vice versa) produces bounded effects.
\end{enumerate}
\end{paradox}

\section{The Triple Attribution Gap}

\begin{theorem}[Three-Party Attribution Inflates]
\label{thm:triple-attribution}
For three ideas $a, b, c$ and observer $d$:
\begin{align*}
&M(a, b \compose c, d) + M(b, a \compose c, d) + M(c, a \compose b, d) \\
&= \operatorname{rs}(a,d) + \operatorname{rs}(b,d) + \operatorname{rs}(c,d) \\
&\quad + \emrg(a, b \compose c, d) + \emrg(b, a \compose c, d) + \emrg(c, a \compose b, d).
\end{align*}
The sum of marginal contributions equals individual resonances PLUS all pairwise emergences.
\end{theorem}

\begin{proof}
Expand each marginal contribution using $M(x, y, d) = \operatorname{rs}(x,d) + \emrg(x,y,d)$:
\begin{align*}
&M(a, b \compose c, d) + M(b, a \compose c, d) + M(c, a \compose b, d) \\
&= [\operatorname{rs}(a,d) + \emrg(a, b \compose c, d)] \\
&\quad + [\operatorname{rs}(b,d) + \emrg(b, a \compose c, d)] \\
&\quad + [\operatorname{rs}(c,d) + \emrg(c, a \compose b, d)].
\end{align*}
Collecting terms gives the result.
\end{proof}

\begin{lstlisting}
/-- DP8. THE TRIPLE ATTRIBUTION GAP -/
theorem triple_attribution_gap (a b c d : I) :
    marginalContrib a (compose b c) d +
    marginalContrib b (compose a c) d +
    marginalContrib c (compose a b) d =
    rs a d + rs b d + rs c d +
    emergence a (compose b c) d +
    emergence b (compose a c) d +
    emergence c (compose a b) d := by
  unfold marginalContrib emergence; ring
\end{lstlisting}

\begin{interpretation}
With three parties, the attribution inflation is even worse than with two. The sum of marginal contributions includes not just the individual resonances but ALL the emergence terms with different coalition partners. As the number of parties grows, the attribution inflation grows combinatorially.

This is why fair division in large organizations is so difficult: the total ``credit'' claimed by all members (their marginal contributions) always exceeds the total value created. The excess equals the sum of all pairwise emergence terms, which grows quadratically in the number of members.
\end{interpretation}


% ================================================================
\chapter{Composition Order and Strategic Advantage}
% ================================================================

\epigraph{The first move is theirs; the last move is mine.}{Bobby Fischer}

\section{The Naive Expectation}

In sequential games, the first mover often has an advantage (the ``first-mover advantage'' in industrial organization: Lieberman and Montgomery, 1988) or a disadvantage (the ``second-mover advantage'' when flexibility is valuable). The direction depends on the specific game.

In ideatic space, the first-mover advantage has a precise algebraic characterization that is independent of the specific ``game'' being played.

\section{The First-Mover Premium Decomposition}

\begin{theorem}[The First-Mover Premium]
The first-mover premium $P(a,b) = [\wt(a \compose b) - \wt(a)] - [\wt(b \compose a) - \wt(b)]$ decomposes as:
\[
P(a,b) = N(a,b) + (\wt(b) - \wt(a)),
\]
where $N(a,b) = \wt(a \compose b) - \wt(b \compose a)$ is the negotiation asymmetry.
\end{theorem}

\begin{proof}
\begin{align*}
P(a,b) &= \wt(a \compose b) - \wt(a) - \wt(b \compose a) + \wt(b) \\
&= [\wt(a \compose b) - \wt(b \compose a)] + [\wt(b) - \wt(a)] \\
&= N(a,b) + (\wt(b) - \wt(a)).
\end{align*}
\end{proof}

\begin{paradox}[The Paradox of the Weak First Mover]
The first-mover premium has two independent components:
\begin{enumerate}
\item The \textbf{negotiation asymmetry} $N(a,b)$: how much more weight $a \compose b$ has compared to $b \compose a$.
\item The \textbf{weight differential} $\wt(b) - \wt(a)$: how much heavier $b$ is compared to $a$.
\end{enumerate}

When negotiation is ``fair'' ($N(a,b) = 0$, e.g., when composition commutes), the first-mover premium is ENTIRELY determined by the weight differential: $P(a,b) = \wt(b) - \wt(a)$.

This means the LIGHTER party benefits more from going first! In a ``fair'' negotiation, the weak benefit more from the first-mover advantage than the strong do.

This inverts the standard first-mover advantage: in most models, the stronger party benefits from moving first because they can commit to a favorable outcome. In ideatic space, when the structural asymmetry is zero, it is the weaker party who benefits from the compositional first-mover advantage.
\end{paradox}

\section{The Order Determines Resonance}

\begin{theorem}[Order Determines Resonance Profile]
The difference in resonance between $a \compose b$ and $b \compose a$, as observed by any $c$, equals the order asymmetry:
\[
\operatorname{rs}(a \compose b, c) - \operatorname{rs}(b \compose a, c) = \emrg(a,b,c) - \emrg(b,a,c).
\]
\end{theorem}

\begin{proof}
Expand using the decomposition:
\begin{align*}
\operatorname{rs}(a \compose b, c) &= \operatorname{rs}(a,c) + \operatorname{rs}(b,c) + \emrg(a,b,c), \\
\operatorname{rs}(b \compose a, c) &= \operatorname{rs}(b,c) + \operatorname{rs}(a,c) + \emrg(b,a,c).
\end{align*}
Subtracting: $\operatorname{rs}(a \compose b, c) - \operatorname{rs}(b \compose a, c) = \emrg(a,b,c) - \emrg(b,a,c)$.
\end{proof}

\begin{interpretation}
The entire difference between two orderings of the same ideas is captured by the emergence asymmetry. The individual resonances $\operatorname{rs}(a,c)$ and $\operatorname{rs}(b,c)$ cancel out. Only the \emph{emergent} contribution differs.

This means that composition order matters ONLY because of emergence. In a world without emergence (all ideas are linear), composition order would not matter at all. The non-commutativity of meaning is ENTIRELY an emergent phenomenon.

This connects to Wittgenstein's insight that ``meaning is use'': the order in which ideas are composed (used) determines their meaning, and this determination operates entirely through emergence. The ``intrinsic'' content of ideas ($\operatorname{rs}(a,c)$ and $\operatorname{rs}(b,c)$) is order-independent; only the emergent content depends on order.
\end{interpretation}

\section{The Zero Asymmetry Theorem}

\begin{theorem}[Zero Asymmetry Implies Same Profile]
If $\emrg(a,b,c) = \emrg(b,a,c)$ for all $c$, then $a \compose b$ and $b \compose a$ have identical resonance profiles: $\operatorname{rs}(a \compose b, c) = \operatorname{rs}(b \compose a, c)$ for all $c$.
\end{theorem}

\begin{proof}
By the previous theorem, the difference is $\emrg(a,b,c) - \emrg(b,a,c)$, which is zero by hypothesis.
\end{proof}

\begin{corollary}
Two ideas have the same composition in both orders (they are ``order-indistinguishable'') if and only if their emergence is symmetric. Order sensitivity is measured entirely by emergence asymmetry.
\end{corollary}


% ================================================================
\chapter{The Emergence Ratchet: Why Complexity Grows}
% ================================================================

\epigraph{Complexity is a sign of sophistication, not of failure.}{Herbert A.\ Simon}

\section{The Naive Expectation}

In many domains, complexity is viewed as a cost: complex systems are harder to maintain, more fragile, and more expensive. Occam's razor counsels simplicity. Information-theoretic minimum description length principles penalize complexity.

The naive expectation: in strategic interaction, complexity should be an optional choice --- agents should be able to stay simple if they prefer.

In ideatic space, complexity growth is not a choice. It is a theorem.

\section{The Cumulative Cultural Evolution Theorem}

\begin{theorem}[Cultural Evolution Is Monotone]
\label{thm:cultural-evolution}
The cumulative gain from $n$ generations of drift is non-negative and non-decreasing:
\begin{enumerate}
\item $\wt(\text{drift}(a,b,n)) - \wt(a) \geq 0$ for all $n$.
\item $\wt(\text{drift}(a,b,n+1)) - \wt(a) \geq \wt(\text{drift}(a,b,n)) - \wt(a)$ for all $n$.
\end{enumerate}
\end{theorem}

\begin{proof}
(1) By induction on $n$. The base case ($n=0$) gives $\wt(a) - \wt(a) = 0$. For the inductive step, $\wt(\text{drift}(a,b,n+1)) \geq \wt(\text{drift}(a,b,n))$ by enrichment, and by the inductive hypothesis $\wt(\text{drift}(a,b,n)) \geq \wt(a)$, so $\wt(\text{drift}(a,b,n+1)) \geq \wt(a)$.

(2) follows from enrichment: $\wt(\text{drift}(a,b,n+1)) \geq \wt(\text{drift}(a,b,n))$, so the cumulative gain at step $n+1$ is at least as large as at step $n$.
\end{proof}

\begin{lstlisting}
/-- EP3. THE DRIFT ACCELERATION THEOREM -/
theorem drift_cumulative_lower_bound (a b : I) (n : N) :
    rs (drift a b n) (drift a b n) >= rs a a := by
  induction n with
  | zero => simp [drift]
  | succ k ih =>
    have hk := drift_enriches a b k
    linarith
\end{lstlisting}

\begin{paradox}[The Impossibility of Cultural Simplification]
Complexity only grows. Every generation adds weight through composition, and no generation can remove it. Cultural artifacts --- languages, laws, technologies, religions --- grow more complex over time because the enrichment axiom guarantees it.

This explains several empirical observations:
\begin{enumerate}
\item \textbf{Legal codes grow}: laws accrete but are rarely simplified.
\item \textbf{Software bloats}: codebases grow but refactoring rarely reduces total complexity.
\item \textbf{Bureaucracies expand}: organizations add rules and procedures but rarely remove them.
\item \textbf{Languages evolve}: languages gain vocabulary and grammatical complexity but rarely simplify.
\end{enumerate}

In each case, the enrichment axiom --- not human psychology or institutional inertia --- provides the fundamental explanation. Complexity growth is a mathematical necessity of compositional systems.
\end{paradox}

\section{The Adaptation Rate Is Always Non-Negative}

\begin{theorem}[Adaptation Is Always Positive]
\label{thm:adaptation-positive}
The adaptation rate $\wt(a \compose b) - \wt(a)$ is always non-negative.
\end{theorem}

\begin{proof}
Direct application of enrichment.
\end{proof}

\begin{interpretation}
Every act of transmission enriches. There is no ``maladaptation'' in the weight sense --- every idea that passes through a cultural context emerges at least as heavy as it entered.

This does not mean that all cultural changes are ``improvements'' in any normative sense. Weight is not quality. A more complex legal code is not necessarily a better one. But it IS heavier, and it CANNOT become lighter through the natural process of cultural transmission.
\end{interpretation}

\section{Connection to Tomasello's Ratchet Effect}

Tomasello (1999) proposed that human cultural evolution exhibits a ``ratchet effect'': innovations are preserved and accumulated across generations, unlike other species where innovations are frequently lost.

In ideatic space, Tomasello's ratchet is a theorem. The enrichment axiom guarantees that cultural innovations (compositions) always increase weight, and the non-degeneracy axiom guarantees that non-void innovations can never be lost (the Extinction Impossibility theorem). Together, these imply that cultural complexity ratchets upward monotonically.

What Tomasello described as a uniquely human cultural capacity is, in ideatic theory, a consequence of the algebraic structure of composition. Any system that composes ideas with enrichment will exhibit the ratchet effect. The ``uniquely human'' part is not the ratchet but the richness of the ideatic space --- the diversity of ideas and emergence that human culture generates.


% ================================================================
\chapter{Grand Synthesis: The Fundamental Theorem of Strategic Interaction}
% ================================================================

\epigraph{Everything should be made as simple as possible, but not simpler.}{Albert Einstein (attributed)}

\section{The Architecture of the Paradoxes}

The preceding fourteen chapters have presented a gallery of counter-intuitive results. Let us now step back and see the architecture that unifies them.

Every paradox in this volume flows from the interaction of two axioms:
\begin{enumerate}
\item \textbf{Enrichment}: $\wt(a \compose b) \geq \wt(a)$.
\item \textbf{Emergence bound}: $\emrg(a,b,c)^2 \leq \wt(a \compose b) \cdot \wt(c)$.
\end{enumerate}

Enrichment gives us the ``positive'' paradoxes: information is irreversible, coalitions always grow, complexity ratchets upward, ideas cannot go extinct.

The emergence bound gives us the ``constraining'' paradoxes: manipulation is expensive, polarization is bounded, attribution is inflationary.

Together, they create the central tension of strategic interaction in ideatic space: \emph{everything grows, but everything is bounded}.

\section{The Grand Synthesis}

\begin{theorem}[The Fundamental Theorem of Strategic Interaction]
\label{thm:fundamental}
For any ideas $a, b$ in an ideatic space:
\begin{enumerate}
\item \textbf{Value creation is guaranteed}: $\wt(a \compose b) \geq \wt(a)$ and $\wt(b \compose a) \geq \wt(b)$.
\item \textbf{Total value exceeds parts}: $\wt(a \compose b) + \wt(b \compose a) \geq \wt(a) + \wt(b)$.
\item \textbf{Surprise is bounded}: $\emrg(a,b,c)^2 \leq \wt(a \compose b) \cdot \wt(c)$.
\item \textbf{Order matters}: $\wt(a \compose b) - \wt(b \compose a) = -[\wt(b \compose a) - \wt(a \compose b)]$.
\item \textbf{Fairness from commutativity}: If $a \compose b = b \compose a$, then $N(a,b) = 0$.
\item \textbf{Attribution inflates}: $M(a,b,c) + M(b,a,c) = \operatorname{rs}(a \compose b, c) + \emrg(b,a,c) \geq \operatorname{rs}(a \compose b, c)$.
\item \textbf{Ideas are immortal}: If $a \neq \void$, then $a \compose b \neq \void$.
\item \textbf{Complexity ratchets}: $\wt(\compN{a}{n+1}) \geq \wt(\compN{a}{n})$.
\end{enumerate}
\end{theorem}

\begin{proof}
Each claim has been proven in the preceding chapters:
\begin{enumerate}
\item The enrichment axiom.
\item Sum of the two enrichment inequalities.
\item The emergence bound axiom.
\item Antisymmetry of subtraction.
\item If $a \compose b = b \compose a$, then $\wt(a \compose b) = \wt(b \compose a)$.
\item The attribution conservation law (Chapter~5) and the fact that emergence can be positive.
\item Non-degeneracy of composition: if the left factor is non-void, the composition is non-void.
\item Enrichment applied to $\compN{a}{n}$ and $a$.
\end{enumerate}
\end{proof}

\section{What the Fundamental Theorem Means}

The Fundamental Theorem says that strategic interaction in ideatic space is simultaneously:
\begin{itemize}
\item \textbf{Productive}: interaction always creates value (properties 1--2).
\item \textbf{Bounded}: the creative surplus is constrained (property 3).
\item \textbf{Asymmetric}: the distribution of value depends on order (properties 4--5).
\item \textbf{Inflationary}: credit always exceeds value (property 6).
\item \textbf{Persistent}: ideas survive and grow (properties 7--8).
\end{itemize}

This combination of properties is unique to ideatic space. In classical game theory, interaction can destroy value (negative-sum games exist), surplus is unbounded (there is no emergence bound), and strategies can die (species go extinct, firms go bankrupt).

The ideatic framework reveals a world where interaction is fundamentally constructive, bounded, and irreversible. This is not a utopian vision --- the asymmetry (property 4) ensures that interaction is often unfair, and the inflationary attribution (property 6) ensures that conflict over credit is inevitable. But it is a world where nihilism is impossible: every interaction creates something, and nothing is ever truly lost.

\section{Connections to the Great Game Theorists}

\subsection*{John Nash}
Nash showed that every finite game has an equilibrium. In ideatic space, the void equilibrium always exists, and every Nash equilibrium has non-negative surplus. Nash's existence theorem is strengthened: not only do equilibria exist, but they are always weakly beneficial.

\subsection*{John von Neumann}
Von Neumann proved the minimax theorem for zero-sum games. In ideatic space, even ``zero-sum'' interactions (orthogonal ideas) create positive weight through enrichment. Von Neumann's zero-sum framework is too restrictive: ideatic games are always positive-sum in weight.

\subsection*{Reinhard Selten}
Selten refined Nash equilibrium to subgame perfection. In ideatic space, the composition order determines the ``subgame structure,'' and the framing effect theorem shows that order matters through emergence chains. Selten's refinement is vindicated: order matters, and it matters for algebraic reasons.

\subsection*{John Harsanyi}
Harsanyi introduced incomplete information into game theory. In ideatic space, ``information'' is compositional: hearing a signal changes your cognitive state irreversibly. Harsanyi's type spaces correspond to ideatic ideas, but with the crucial addition of enrichment: learning always increases weight.

\subsection*{Robert Aumann}
Aumann proved that rational agents with common priors cannot agree to disagree. In ideatic space, ``agreement'' is measured by resonance, and the polarization bound constrains how much agents can disagree. Aumann's result is extended: disagreement is not just unlikely under common priors --- it is \emph{bounded} by the emergence constraint.

\subsection*{Roger Myerson}
Myerson developed the revelation principle and optimal mechanism design. In ideatic space, the manipulation bound provides a \emph{physical} constraint on mechanism design that goes beyond incentive compatibility: even a mechanism designer with full power cannot create emergence beyond the bound. Myerson's revelation principle still holds, but the achievable outcomes are further constrained by the algebraic structure of composition.

\section{Open Questions}

We conclude with several questions that this volume raises but does not resolve:

\begin{enumerate}
\item \textbf{Tight bounds}: The emergence bound is $\emrg(a,b,c)^2 \leq \wt(a \compose b) \cdot \wt(c)$. Is this bound tight? Are there ideatic spaces where it is achieved with equality?

\item \textbf{Convergence}: The weight ratchet shows that $\wt(\compN{a}{n})$ is non-decreasing. Does it converge? If so, to what? This would be a ``fixed point'' of the ratchet.

\item \textbf{Non-commutative dynamics}: The framing effect theorem shows that order matters. Can we characterize which orderings are ``optimal'' (maximize total weight) for a given set of ideas?

\item \textbf{Large coalitions}: The enrichment chain shows that adding members always enriches. But does the marginal enrichment decay? Is there a ``law of diminishing marginal enrichment'' for large coalitions?

\item \textbf{Information-theoretic capacity}: Is there an ideatic analog of Shannon capacity --- a maximum rate at which ideas can be composed while maintaining some measure of ``fidelity''?

\item \textbf{Evolutionary dynamics}: The complexity ratchet shows that cultural evolution is monotone in weight. But weight is not the only measure of ``fitness.'' Can we characterize evolutionary dynamics in ideatic space using richer fitness measures?
\end{enumerate}

These questions suggest that the theory of strategic interaction in ideatic space is far from complete. The paradoxes of this volume are the first glimpses of a rich mathematical landscape that remains to be explored.

\bigskip

\begin{center}
\rule{0.5\textwidth}{0.5pt}
\end{center}

\bigskip

\noindent\textit{Every theorem in this volume has been verified in Lean~4 with zero sorry axioms. The source files are:}
\begin{itemize}
\item \texttt{IDT\_v3\_Foundations.lean} --- core axioms and definitions.
\item \texttt{IDT\_v3\_SignalGames.lean} --- signal games and communication paradoxes.
\item \texttt{IDT\_v3\_DeepGames.lean} --- deep game-theoretic results and attribution.
\item \texttt{IDT\_v3\_Cooperation.lean} --- cooperative game theory and coalition paradoxes.
\item \texttt{IDT\_v3\_Networks.lean} --- network dynamics and cascade theorems.
\item \texttt{IDT\_v3\_Evolution.lean} --- evolutionary paradoxes and complexity ratchet.
\end{itemize}


\end{document}
